\documentclass[11pt]{article}
\usepackage[letterpaper,margin=1in]{geometry}
\usepackage{graphicx}
\usepackage{booktabs}
\usepackage{amsmath}
\usepackage{xcolor}
\usepackage[hidelinks]{hyperref}
\usepackage[capitalise,nameinlink]{cleveref}
\usepackage{caption}
\captionsetup{font=small,labelfont=bf}
\graphicspath{{../results/figures_nocap/}}

% Every number below carries a "% src:" comment naming the table it comes from, under
% results/tables/ unless stated. Regenerate the tables with analysis/run_all.py.

\newcommand{\pmrep}{\texttt{.pm} replay}

\title{Evaluating Whether Vernier-Driven Pattern Replay Predicts\\ABCG's Choice of Halo-Exchange Strategy on Dane}
\author{Nicholas Bacon\\University of New Mexico}
\date{October 2026}

\begin{document}
\maketitle

\begin{abstract}
Choosing a halo-exchange strategy once, at program launch, pays off only if a cheap
measurement predicts which strategy the application itself would choose. However, to
our knowledge no one has measured whether a Vernier characterization of an
application, replayed through a communication benchmark, makes the same choice as the
application. In this paper we measure how well four stand-ins choose among three
halo-exchange strategies (plain MPI point-to-point, MPI Advance neighbor alltoallv, and
MPI Advance locality-aware aggregation) for ABCG's conjugate gradient (CG) iteration.
The four stand-ins are a replay of Vernier's aggregate histograms, a replay of Vernier's
per-rank capture, a replay of the exact matrix halo, and ABCG's own SpMV exchange. We
ran 40 separate allocations at each of 2, 4, 8, and 16 nodes of LLNL's Dane system on
six SuiteSparse matrices. ABCG's SpMV exchange chooses an arm within 4.2\% of CG's best
at all 24 combinations of matrix and node count. The exact-matrix replay and the per-rank
Vernier replay stay within 10\% at 23 and 22 of the 24; every choice that costs CG more
than 5\% ranks the aggregating strategy too low, and the exact-matrix replay misses
exactly where CG's gain from aggregation is at most 1.21 times. The aggregate-histogram
replay chooses differently from the per-rank replay at 4 of the 24 combinations, all at
2 and 4 nodes, where CG takes up to 32.3 times as long per iteration on its choice.
Every stand-in prefers MPI Advance's non-aggregating alltoallv where CG is fastest on
plain MPI, at a cost of at most 4.2\%. Finally, 266 of 9,600 MPI and STANDARD runs lost
whole multiples of 0.160\,s inside the timed loop; these stalls decide CG's arm ordering
on the crankseg matrices at 2 nodes, and the aggregate-histogram replay almost never
shows them. We have not identified their cause.
\end{abstract}

\section{Introduction}

No single halo-exchange strategy is fastest for every sparse workload. Collom, Li, and
Bienz showed, at operator level, that the best of direct point-to-point exchange and
locality-aware aggregation changes across the levels of one algebraic multigrid
hierarchy~\cite{collom2023}. The research plan this work belongs to proposes a decision
framework that chooses a strategy once per program launch from a small vector of
cheaply measured pattern features, and judges that framework by its worst-case regret,
the largest slowdown a wrong choice causes~\cite{researchplan}. Vernier records such
features at scale as histograms of message size and logical neighbor distance, keyed by
neighbor count, kept per rank or merged across ranks~\cite{vernier2025}. The
irregular-benchmarking replay benchmark regenerates a communication pattern from a
Vernier capture and times it under different strategies.

However, a Vernier capture is a lossy summary of the pattern. The aggregate data manager
keeps one histogram set for the whole job, and even the per-rank hierarchical manager
keeps counts and sizes but not which matrix entries move. Whether a replay built from
either capture ranks the strategies the way the application does is the plan's
hypothesis H17, and whether the cheaper aggregate capture suffices is H18; to our
knowledge neither had been measured. The plan also predicts, before any data, that
replay will order the non-aggregating strategies correctly and degrade on the
aggregating one, because aggregation depends on the mapping of ranks to NUMA domains and
nodes, which a histogram of logical neighbor distance does not record.

This paper measures, for ABCG's CG iteration on Dane, how much CG slows down when each
of four stand-ins chooses the strategy instead of CG itself. By \emph{stand-in} we mean
a cheaper measurement that exchanges a version of the SpMV halo of the same matrix at
the same node count; each stand-in in this study is one step further from CG than the
one before it. The stand-ins do not reproduce the allreduces or the computation in CG.
This paper makes the following contributions:
\begin{itemize}
\item The identification of a timing anomaly on Dane, in which some MPI and STANDARD
  runs lose whole multiples of 0.160\,s inside the timed loop, and an account of which
  stand-ins show it (\cref{sec:stalls});
\item A step-by-step comparison of what each stand-in changes about the pattern and
  about the time per iteration (\cref{sec:steps});
\item An evaluation of each stand-in's strategy choice against CG's own choice by
  decision regret over 40 allocations per configuration (\cref{sec:regret,sec:matters});
\item A set of bounded conclusions about when the replay benchmark can stand in for
  ABCG's CG (\cref{sec:discussion}).
\end{itemize}
\Cref{sec:standins} describes the stand-ins and the strategies, \cref{sec:setup} the
experimental setup, \cref{sec:limits} what this data cannot show, and
\cref{sec:conclusion} concludes.

\section{The Stand-Ins and the Strategies}
\label{sec:standins}

We compare five ways of exchanging the SpMV halo of a matrix, ordered from the run we
care about to the cheapest stand-in. All five exchange the halo of the same matrix,
partitioned into one contiguous block of rows per MPI rank.

\paragraph{ABCG CG (the reference):} ABCG's \texttt{benchmark\_cg\_pilot} runs CG
iterations~\cite{abcg}, each consisting of an SpMV with its halo exchange, two
allreduces, and vector updates, with no barrier between iterations. We treat its time
per iteration as the true run. On all six matrices CG reached ABCG's 500-iteration cap
without converging, so these are fixed-iteration timings, not time to solution.
% src: audit.csv cg_runs_reference_iters_500 = cg_runs at every node count

\paragraph{ABCG SpMV exchange:} ABCG's \texttt{benchmark\_spmv} runs CG's
pack-and-exchange code followed by a barrier, without the local multiply, the
allreduces, or the vector updates.

\paragraph{\pmrep:} irregular-benchmarking reads the matrix, builds the same halo, and
exchanges it through Cabana's gather, with a barrier after every exchange and no
computation. In all 2,880 paired runs, the minimum, maximum, and average over ranks of
its messages and bytes sent equal ABCG's exactly; it differs from the ABCG SpMV exchange
in code path. For STANDARD and LOCALITY it calls MPI Advance's persistent-request
(\texttt{\_INIT}) neighbor alltoallv; its run headers also report
``\texttt{persistent: false}'', a setting whose meaning in the benchmark we have not
confirmed.
% src: results/runs.csv, pm vs abcgspmv numberOfmessages/data_sent min/max/avg, 0 mismatches in 2,880 pairs; run headers "-alltoallv: NEIGHBOR_ALLTOALLV_INIT_*", "-persistent: false"

\paragraph{Fine replay:} irregular-benchmarking reads a Vernier \emph{hierarchical}
capture, which records for each rank the number of elements sent to each partner. The
capture comes from one exchange of ABCG's SpMV at the same node count. Each rank sends
its recorded number of elements to its recorded partners; the minimum, maximum, and
average over ranks of its messages, bytes, on- and off-node bytes, and same- and
other-NUMA bytes equal the \pmrep's in every run. It does not send the same matrix
entries, so its memory access pattern differs from the \pmrep's.
% src: audit.csv fine_vs_pm_mismatched_values = 0 at every node count

\paragraph{Coarse replay:} irregular-benchmarking reads Vernier's \emph{aggregate}
histograms, one histogram set for the whole job. Each rank samples a neighbor count,
distinct partners by logical neighbor distance, and message sizes within histogram bins
from these job-wide histograms, and draws a new pattern in every run. It does not
record which rank sends what or which ranks share a NUMA domain.

\paragraph{Strategies:} Each stand-in runs three strategies, which we call
\emph{arms}. \textbf{MPI} is plain MPI point-to-point. \textbf{STANDARD} is MPI
Advance's neighbor alltoallv without aggregation~\cite{mpiadvance2023}.
\textbf{LOCALITY} is MPI Advance's locality-aware neighbor alltoallv, which aggregates
messages through NUMA-region leaders (\texttt{-s NUMA}); it is the only aggregating arm.
ABCG uses persistent requests for STANDARD and LOCALITY. The campaign also ran ABCG with
one non-persistent MPI Advance call per exchange; its LOCALITY exchange takes 0.31 to 54
times as long as the persistent one, so we treat it as a different strategy and leave it
out of this paper.
% src: ratios.csv stand_in=abcgspmv reference=abcgspmvnp arm=LOCALITY slow_runs=kept, 1/max median_ratio = 0.31, 1/min = 53.7

\section{Experimental Setup}
\label{sec:setup}

\paragraph{Machine and software:} We ran on Dane at LLNL with 112 MPI ranks per node,
each bound to its own core, at 2, 4, 8, and 16 nodes (224 to 1,792 ranks), in the
shared \texttt{pbatch} queue. We excluded nodes dane261 and dane764, which report 2 NUMA
domains instead of 8 and crash LOCALITY. The software was Open MPI 4.1.2, gcc 13.3.1,
MPI Advance \texttt{locality\_aware} commit 752a19b, irregular-benchmarking 2593e29,
ABCG 29953fe, and Vernier 09d650b.
% src: provenance_*.txt (distinct CPU sets: 112 per node), data/README.txt

\paragraph{Matrices:} We used six SuiteSparse matrices: crankseg\_2 (63,838 rows,
14.1M nonzeros), nd12k (36,000 rows, 14.2M nonzeros), x104 (108,384 rows, 10.2M
nonzeros), crankseg\_1 (52,804 rows, 10.6M nonzeros), m\_t1 (97,578 rows, 9.8M
nonzeros), and nd24k (72,000 rows, 28.7M nonzeros).
% src: run file headers ("-matrix:" line)

\paragraph{Procedure:} For each node count, a pilot job first captured each matrix's
SpMV halo with Vernier, then ran every kind on every matrix on STANDARD to choose an
iteration count $n$ that takes about 10\,s; all three arms use that $n$. Each run warms
up connections to ranks $r\pm1$ without timing, performs $k=25$ startup iterations,
which $q$ excludes, and then performs $n$ timed iterations. Forty launches per node
count, each a separate Slurm allocation, ran every combination of 7 kinds (the four
stand-ins, CG, and the two non-persistent ABCG kinds), 6 matrices, and 3 arms, 126 runs
per launch, each launch in its own random order with one \texttt{srun} per run. At most
four launches of one node count ran at once, and up to 16 overall; the 160 launches ran
over 23.7 hours. Every run exited normally, every benchmark MPI Advance run verified its
ghost values against plain MPI, and no ABCG SpMV run reported a wrong halo value.
% src: audit.csv (exit_nonzero, verify_fail, bench_mpia_runs_without_pass, abcg_check_failed all 0); runs.csv (start/end: span 23.7 h, max 16 concurrent launches)

\paragraph{Metric:} The time per iteration $q$ of a run is its total timed time,
maximized over ranks, divided by $n$. Because each timed iteration ends at a barrier or
an allreduce, the minimum and maximum over ranks differ by at most 0.74\%.
% src: audit.csv rank_spread_timedTotal_max max over node counts = 0.007373

\paragraph{Statistics:} The unit of replication is the launch, because runs within a
launch share an allocation. For a stand-in $s$ and a reference $r$ at one matrix and
node count, the stand-in chooses the arm with the lowest median $q_s$ over launches,
and its \emph{decision regret} is
\begin{equation}
  \mathrm{regret} = \frac{\mathrm{median}\, q_r(\text{stand-in's arm})}{\min_a \mathrm{median}\, q_r(a)} - 1 .
\end{equation}
A regret of 0 means the stand-in chose the reference's fastest arm, and a regret of
0.21 means the reference takes 21\% longer per iteration on the stand-in's choice. This
is the launch-time setting of the plan, in which one choice is made per program and
configuration. Intervals come from resampling launches with replacement (10,000
resamples), which keeps each launch's runs together. To summarize the results we call
a stand-in \emph{useful for choosing an arm for CG} where its decision regret on CG is
at most 10\% at every combination of matrix and node count. The 10\% figure is the
equivalence margin we fixed before computing any result; applying it to decision regret
is a choice we made after seeing the regrets.

\section{Some MPI and STANDARD Runs Lose Whole 0.160-Second Stalls}
\label{sec:stalls}

\begin{figure}[t]
\centering
\includegraphics[width=0.78\linewidth]{stall_quantum_excess_fine_crankseg_2_STANDARD_16nodes.png}
\caption{Extra timed time of each of 40 launches of the fine replay of crankseg\_2 on
STANDARD at 16 nodes ($n=57$ timed iterations), relative to the fastest launch, in
units of 0.160\,s. Every bar falls on a whole number.}
\label{fig:stallhist}
\end{figure}

\begin{figure}[t]
\centering
\includegraphics[width=0.78\linewidth]{stall_quantum_coherence_by_cell.png}
\caption{Phase coherence of the launch-to-launch differences in total timed time with a
0.160\,s period (1 = every difference is a multiple of 0.160\,s; about 0.1 = no
pattern), for each of the 447 combinations of node count, kind, matrix, and arm whose
total varies by at least 0.3\,s across its 40 launches. Points above the dashed line
(0.5) are almost all MPI or STANDARD.}
\label{fig:stallcoh}
\end{figure}

We report this effect first because it shapes the crankseg and 2-node results below.
For each combination of node count, kind, matrix, and arm whose total timed time varies
by at least 0.3\,s across launches (447 of them), we measured the phase coherence of
the pairwise launch-to-launch differences against candidate periods from 0.08 to
0.30\,s. In 61 combinations the best coherence exceeds 0.5 (from 0.53 up to 0.995);
their best periods have a median of 0.1600\,s and an interquartile range of 0.1600 to
0.1605\,s (\cref{fig:stallhist,fig:stallcoh}). In these combinations, slow runs lose a
whole number of 0.160\,s stalls rather than running uniformly slower; in the worst
launches of \cref{fig:stallhist}, 56 stalls fall within 57 timed iterations. 22 of the
61 combinations are of the non-persistent ABCG kinds.
% src: stall_quantum.csv (447 rows; 61 with coherence > 0.5, min 0.53; best_period_s median and IQR; max coherence 0.9952; 22 of 61 abcg*np)

We flagged as slow every run whose $q$ exceeds five times the median of its
combination. Among the kinds this paper studies, 266 of 9,600 MPI and STANDARD runs are
flagged: 76 of the \pmrep, 98 of the fine replay, 90 of ABCG's SpMV exchange and CG, and
2 of the coarse replay. No LOCALITY run of these kinds is flagged. The slow runs fall on
crankseg\_1 and crankseg\_2 at every node count, on nd24k at 4 to 16 nodes, and on nd12k
at 8 and 16 nodes. The \pmrep{} and the fine replay have slow runs in 18 combinations of
node count, matrix, and arm, and the persistent ABCG kinds in 16, of which 11 are
shared. The stalls also set $n$. At 16 nodes, the pilot's STANDARD run of the fine
replay of crankseg\_2 measured 142.8\,ms per iteration, so the pilot chose $n=57$. The
fastest MPI run of that pattern takes 0.41\,ms per iteration, so its timed loop lasts
about 23\,ms rather than 10\,s, and each stall adds 2.8\,ms to $q$, about seven times
the stall-free time; even the fastest STANDARD run (3.2\,ms) contains one stall. Across
the campaign, the pilot's time per iteration was more than 2 times the fastest later
STANDARD run's for 47 of the 168 combinations of node count and pattern. On 4-node crankseg\_2, the median MPI and STANDARD $q$ of the \pmrep,
the fine replay, and ABCG's SpMV exchange is 4 to 23 times the fastest launch's.
% src: slow_runs.csv / runs.csv slow flag (studied kinds, MPI+STANDARD: 266 of 9,600; pm 76, fine 98, abcgspmv+abcgcg 90, coarse 2; 0 LOCALITY); cell overlap 18/16/11 computed from runs.csv; pilot_01_crankseg_2 file in 8013862-pilot-16

Neither launch start time (Spearman $|\rho| \le 0.12$, $p \ge 0.45$ at every node
count) nor membership of any one node (124 nodes tested, smallest Bonferroni-corrected
$p = 0.73$) predicts which launches stall. The kind of the preceding run shows no
consistent association with stalling. We do not know the cause. The five-times-median
rule used in the robustness checks below does not remove the stalls in combinations
where they inflate the median itself.
% src: launch_stall_score.csv; slow_by_node.csv; slow_by_predecessor.csv

\section{What Each Step Changes}
\label{sec:steps}

\begin{figure}[t]
\centering
\includegraphics[width=0.9\linewidth]{coarse_structure_ratio_4nodes.png}
\caption{Per-rank pattern quantities of the coarse replay divided by the exact values
(which the fine replay matches), at 4 nodes. Markers are medians over 40 launches and
bars span all 40, because the coarse replay draws a new pattern in every run.}
\label{fig:structure}
\end{figure}

\begin{figure}[t]
\centering
\includegraphics[width=0.78\linewidth]{ladder_time_ratio_vs_cg.png}
\caption{Time per iteration of each stand-in divided by ABCG CG's on the same arm. Each
point is one combination of matrix and node count (median over 40 launches of the
per-launch ratio).}
\label{fig:timeratio}
\end{figure}

\paragraph{Coarse to fine replay:} The coarse replay keeps the mean bytes per rank
within 0.7\% below to 3.1\% above the exact value and the mean message count within
1.0\%. On crankseg\_1, crankseg\_2, and nd24k at 2 to 16 nodes (medians over launches;
\cref{fig:structure} shows 4 nodes), however, it sends 0.54 to 0.69 times the exact
same-NUMA bytes, 0.73 to 0.91 times the on-node bytes, and 1.05 to 1.83 times the
off-node bytes. On every matrix, its busiest rank sends 1.17 to 2.06 times as many bytes
as the exact busiest rank. The loss of same-NUMA traffic is consistent with sampling
partners by logical neighbor distance, which does not record which ranks share a NUMA
domain; we did not test this explanation. The coarse replay's time per iteration is 0.59
to 1.17 times the fine replay's on LOCALITY and 0.018 to 1.39 times on MPI and STANDARD.
The low end on MPI and STANDARD comes from combinations in which the fine replay stalls
and the coarse replay does not.
% src: h18_structure.csv (ratio_median ranges over all node counts); ladder_ratio.csv relation=vs_next stand_in=coarse

\paragraph{Fine replay to \pmrep:} On LOCALITY their times agree within 1\% at all 24
combinations (median per-launch ratio 0.993 to 1.008, equivalent within $\pm 10\%$ at
all 24). On MPI and STANDARD they are equivalent at 15 of 24 each. The other 18 are not
shown equivalent (median ratios 0.86 to 1.42), and 13 of them contain a slow run of the
fine replay or the \pmrep.
% src: ratios.csv stand_in=fine reference=pm slow_runs=kept; slow-run membership from runs.csv

\paragraph{\pmrep{} to ABCG SpMV exchange:} The two send the same messages and bytes
(by the minimum, maximum, and average over ranks), so the difference is the code path.
On LOCALITY the \pmrep{} is 1.49 to 1.99 times slower than ABCG's exchange at every
combination. The replay calls the locality-aware exchange through Cabana's gather and
ABCG calls it on its own buffers; we have not isolated which part of the Cabana path
costs the time. On MPI and STANDARD the ratio is 0.95 to 3.79, equivalent within
$\pm 10\%$ at 10 and 11 of 24 combinations.
% src: ratios.csv stand_in=pm reference=abcgspmv slow_runs=kept (LOCALITY 1.488-1.992, 0/24 equivalent; MPI 0.976-3.192, 10/24; STANDARD 0.950-3.786, 11/24)

\paragraph{ABCG SpMV exchange to CG:} A CG iteration adds the local multiply, two
allreduces, and vector updates to the same pack-and-exchange code, and drops the
barrier. ABCG's standalone SpMV exchange takes 0.62 to 0.75 times as long as a CG
iteration on LOCALITY, and 0.39 to 0.97 times as long on MPI and STANDARD
(\cref{fig:timeratio}).
% src: ladder_ratio.csv relation=vs_cg stand_in=abcgspmv

No stand-in's time per iteration is a fixed multiple of CG's. Across the 24
combinations, the \pmrep's time on MPI ranges from 0.44 to 3.26 times CG's, and the
coarse replay's from 0.035 to 0.80 times (\cref{fig:timeratio}). A stand-in can
therefore support a launch-time choice only by preserving the order of the arms, which
the next section measures.
% src: ladder_ratio.csv relation=vs_cg (pm MPI 0.443-3.260; coarse MPI 0.035-0.798)

\section{Whether Each Stand-In Chooses CG's Arm}
\label{sec:regret}

\begin{figure}[t]
\centering
\includegraphics[width=0.78\linewidth]{ladder_worst_regret_vs_cg.png}
\caption{Worst-case decision regret of each stand-in on ABCG CG over the 24
combinations of matrix and node count, and the number of combinations in which the
regret exceeds 10\% (dashed line). Note the logarithmic axis.}
\label{fig:worst}
\end{figure}

\begin{figure}[t]
\centering
\includegraphics[width=\linewidth]{ladder_regret_vs_cg.png}
\caption{Decision regret of each stand-in on ABCG CG at each combination of node count
and matrix, with 95\% launch-bootstrap intervals. 40 launches per combination.}
\label{fig:regretpoints}
\end{figure}

\begin{table}[t]
\centering
\caption{Worst-case decision regret over 24 combinations, scored on CG and on the next
step toward CG, with 95\% launch-bootstrap intervals. Slow runs kept.}
\label{tab:ladder}
\small
\setlength{\tabcolsep}{4pt}
\begin{tabular}{lllcc}
\toprule
Stand-in & Scored on & Worst case at & Worst-case regret (95\% CI) & Over 10\% \\
\midrule
ABCG SpMV & CG & 2 nodes, x104 & 4.2\% (3.1--5.1\%) & 0 of 24 \\
\pmrep & CG & 4 nodes, nd12k & 21\% (19--26\%) & 1 of 24 \\
Fine replay & CG & 2 nodes, crankseg\_1 & 101\% (0--133\%) & 2 of 24 \\
Coarse replay & CG & 4 nodes, crankseg\_2 & 3,129\% (1,738--7,357\%) & 6 of 24 \\
\midrule
\pmrep & ABCG SpMV & 4 nodes, nd12k & 30\% (27--33\%) & 1 of 24 \\
Fine replay & \pmrep & 2 nodes, crankseg\_1 & 1.4\% (0--35\%) & 0 of 24 \\
Coarse replay & Fine replay & 4 nodes, crankseg\_2 & 1,720\% (1,535--4,565\%) & 4 of 24 \\
\bottomrule
\end{tabular}
% src: ladder_regret.csv slow_runs=kept, max decision_regret per (stand_in, relation) and count > 0.10
\end{table}

Only ABCG's SpMV exchange is useful for choosing an arm for CG by our criterion; its
choice costs CG at most 4.2\% at any combination (\cref{fig:worst}, \cref{tab:ladder}).
The \pmrep{} and the fine replay exceed 10\% at one and two combinations, and the coarse
replay at six. Dropping the runs flagged as slow changes none of these counts; the
coarse replay's worst case drops to 1,846\%.
% src: ladder_regret.csv slow_runs=dropped (worst coarse 18.464; over-10% counts unchanged: 0, 1, 2, 6)

\paragraph{Misses involving LOCALITY:} Every choice that costs CG more than 5\%, 12 in
all, is a stand-in choosing STANDARD where CG is fastest on LOCALITY: 7 by the coarse
replay, 3 by the fine replay, and 2 by the \pmrep. No stand-in chose LOCALITY where CG
was fastest without it. These include all three replays at 16-node m\_t1, where CG's
LOCALITY advantage is 7.1\%. This is the half of the plan's prediction that held. The
plan's stated reason, a missing rank-to-region mapping, does not explain the
\pmrep's misses, because the \pmrep{} sends the exact pattern.
% src: ladder_regret.csv relation=vs_cg slow_runs=kept, decision_regret > 0.05 (12 rows, all STANDARD vs LOCALITY; 16-node m_t1 0.071 for pm, fine, coarse)

\paragraph{The \pmrep:} The \pmrep{} misses at exactly the three combinations where
CG's LOCALITY advantage is smallest: 1.016, 1.071, and 1.213 times at 16-node x104,
16-node m\_t1, and 4-node nd12k. It chooses LOCALITY at all 13 combinations where CG's
advantage is 2.0 times or more; no combination falls between 1.21 and 2.0 times. At
4-node nd12k, the \pmrep's LOCALITY takes 1.58 times as long as ABCG's while its MPI and
STANDARD match ABCG's within 3.1\%, which is enough to overturn LOCALITY's 21\% margin in
CG. On crankseg, part of the \pmrep's correct choice comes from its own stalls. At
4-node crankseg\_2 its fastest STANDARD launch (0.260\,ms) beats its fastest LOCALITY
launch (0.365\,ms), and it chooses LOCALITY by median only because its median STANDARD
time is 23 times its fastest.
% src: h0_best_arm.csv reference=abcgcg (margin = min(MPI,STANDARD)/LOCALITY medians) joined with ladder_regret.csv pm choices; ratios.csv pm/abcgspmv 4 nd12k (MPI 1.031, STANDARD 1.014, LOCALITY 1.581); runs.csv 4-node crankseg_2 pm min/median q

\paragraph{The fine replay:} The fine replay's choice differs from the \pmrep's only at
2-node crankseg\_1. There the \pmrep's STANDARD and LOCALITY medians are within 1.4\% of
each other, so the two replays land on opposite sides of a near-tie: the \pmrep's
choice flips in 43\% of bootstrap resamples and the fine replay's in 23\%. CG's cost of
the fine replay's choice there is 101\% (95\% CI 0 to 133\%).
% src: ladder_regret.csv fine/pm vs_cg 2 crankseg_1 (p_choice_changes 0.229 and 0.4285); fine vs_next 0.0144

\paragraph{The coarse replay:} The coarse replay chooses a different arm from the fine
replay at 4 of 24 combinations: crankseg\_2 at 2 nodes, and crankseg\_1, crankseg\_2,
and nd24k at 4 nodes. At 4-node crankseg\_2, CG takes 32.3 times as long per iteration
on its choice (STANDARD) as on LOCALITY. Its two other misses above 10\% on CG, 2-node
crankseg\_1 and 4-node nd12k, are shared with the fine replay. The coarse and fine
replays agree at all 12 combinations at 8 and 16 nodes. The 4-node crankseg\_2 miss is
not a stall artifact, because CG's fastest STANDARD launch (0.951\,ms) is still 3.3
times its fastest LOCALITY launch (0.288\,ms). By contrast, at 2 nodes, CG's arms scored
on their fastest 10\% of launches make MPI fastest on crankseg\_1 and crankseg\_2, and
the 2-node misses of the coarse and fine replays shrink to at most 2.6\%. This
fastest-10\% analysis was not planned before we saw the stalls.
% src: ladder_regret.csv coarse vs_next; h18_coarse_vs_fine.csv nodes >= 8 all 0; runs.csv 4-node crankseg_2 abcgcg min q; h17_p10.csv reference=abcgcg 2-node crankseg (<= 0.026)

\paragraph{Misses between MPI and STANDARD:} The other half of the plan's prediction,
that replay orders the non-aggregating arms correctly, failed. CG is fastest on MPI at 8
combinations (2-node nd12k, x104, m\_t1, and nd24k, and 4- and 8-node x104 and m\_t1),
and at all 8 every stand-in, including ABCG's SpMV exchange, chooses STANDARD. Within
those combinations the \pmrep{} puts STANDARD ahead of MPI in 90 to 100\% of launches,
the fine replay in 82 to 100\%, and ABCG's SpMV exchange in 50 to 88\%, while CG puts
STANDARD ahead in 0 to 60\%. None of the 8 combinations contains a slow run, so stalls
do not explain this bias, and we have not identified its cause. It costs CG only 0.1 to
4.2\%, because MPI and STANDARD are close in CG at these combinations.
% src: ladder_regret.csv vs_cg at the 8 MPI-best combos (all STANDARD); per-launch STANDARD<MPI shares from runs.csv; no slow runs of studied kinds in those combos; h0_best_arm.csv MPI/STANDARD gap 0.001-0.042

\section{Whether the Choice Matters}
\label{sec:matters}

\begin{figure}[t]
\centering
\includegraphics[width=0.85\linewidth]{h0_selection_gain_abcgcg.png}
\caption{How much faster ABCG CG runs when each workload uses its own fastest arm than
when every workload uses the single best arm (geometric mean over workloads). The arms
are chosen on half of the launches and scored on the other half; bars are 95\% ranges
over 2,000 bootstrap splits.}
\label{fig:gain}
\end{figure}

\begin{figure}[t]
\centering
\includegraphics[width=0.85\linewidth]{h8_nstar_abcgcg_LOCALITY_vs_MPI.png}
\caption{Number of CG iterations after which LOCALITY's extra setup cost is repaid,
against MPI, from medians over 40 launches with 95\% launch-bootstrap intervals. Red
markers at the top mean LOCALITY never overtakes MPI. Dotted lines mark the 22 and 40
iterations reported by Collom et al.~\cite{collom2023}.}
\label{fig:nstar}
\end{figure}

A stand-in is worth having only where CG's best arm changes between configurations. By
median, CG is fastest on LOCALITY for crankseg\_1 and crankseg\_2 at every node count,
for nd12k and nd24k from 4 nodes on, and for m\_t1 and x104 only at 16 nodes; everywhere
else it is fastest on MPI. Choosing the arm per matrix instead of using the single best
arm makes CG 28\% (95\% range 8 to 46\%) faster at 2 nodes, 29\% (29.0 to 29.8\%) at 4
nodes, and 11\% (10 to 12\%) at 8 nodes. At 16 nodes LOCALITY is fastest for all six
matrices and the gain is 0 (\cref{fig:gain}). Scored on each arm's fastest 10\% of
launches, CG's best arm is unchanged at every 4- and 8-node combination. At 2 nodes it
is MPI or STANDARD for every matrix, so the 2-node gain comes from stalls. The data
therefore supports the plan's H0, that the best arm changes across workloads, at 4 and 8
nodes, but not at 2 or 16.
% src: h0_best_arm.csv reference=abcgcg slow_runs=kept; h0_selection_gain.csv reference=abcgcg across=matrices slow_runs=kept; h17_p10.csv reference columns for abcgcg

A launch-time choice of LOCALITY also has to repay LOCALITY's setup. Against MPI, CG
repays it after 7 to 31 iterations on the crankseg matrices at 4 to 16 nodes, and after
86 and 103 iterations at 2 nodes. This is the same order as the 22 and 40 iterations
Collom et al.\ reported at 2,048 cores on a different machine and MPI
library~\cite{collom2023}. It takes 26 to 83 iterations on nd24k and 70 to 778 on nd12k
at 4 to 16 nodes, and LOCALITY never overtakes MPI on either at 2 nodes. On m\_t1 and
x104, LOCALITY never overtakes MPI at 2 to 8 nodes and needs 2,449 (95\% CI 1,968 to
5,599) and 8,518 (3,389 to 65,722) iterations at 16 nodes (\cref{fig:nstar}). These
counts exclude the untimed warmup with ranks $r\pm1$.
% src: h8.csv kind=abcgcg versus=MPI slow_runs=kept

\section{Discussion}
\label{sec:discussion}

The results above lead us to the following conclusions about when the replay benchmark
can stand in for ABCG's CG on Dane.

\paragraph{1) Whether to aggregate:} The \pmrep{} chose LOCALITY at every combination
where CG's advantage from LOCALITY was 2.0 times or more, and missed at all three where
it was 1.02 to 1.21 times. Applying this requires knowing CG's margin, which is what the
stand-in is meant to replace, so in practice it means the benchmark cannot be trusted to
recommend LOCALITY when aggregation's benefit to CG is below some threshold between 1.21
and 2.0 times. The fine replay missed once even above that range, at a near-tie in the
replays. On crankseg, the replays' correct choices of LOCALITY rest partly on their own
stalls; the replays and the persistent ABCG kinds stall in 11 shared combinations of the 23 in which either stalls.

\paragraph{2) MPI versus STANDARD:} Every stand-in, including ABCG's own SpMV exchange,
prefers STANDARD where CG is fastest on MPI. The bias is systematic, not noise, and its
cause is unknown. It cost CG 0.1 to 4.2\% here; it would matter more on a workload where
MPI and STANDARD differ more in CG.

\paragraph{3) The per-rank capture against the matrix:} The fine replay's choice matches
the \pmrep's at 23 of 24 combinations, and the one difference is a near-tie in both
replays. Where the matrix is not available, the hierarchical capture is as good an input
as the matrix for this decision at 224 to 1,792 ranks. However, it stores one file per
rank (1,792 files at 16 nodes), so its input cost grows with process count.
% src: ladder_regret.csv fine vs_next (choice differs only at 2-node crankseg_1); data/README.txt (112 x nodes hierarchical files)

\paragraph{4) The aggregate capture at 2 and 4 nodes:} The coarse replay chooses
differently from the fine replay at 4 of 24 combinations, all at 2 and 4 nodes, at a
cost to CG of up to 32.3 times the time per iteration. It undercounts same-NUMA bytes and
overloads the busiest rank by similar factors at 8 and 16 nodes, where it agrees with the
fine replay, so these distortions do not by themselves explain the disagreements. It also
almost never stalls. For this decision on Dane, H18 fails for the aggregate manager at 2
and 4 nodes; the plan states H18 for the combined manager, which this campaign did not
test.

\paragraph{5) Time per iteration:} No stand-in's time is a fixed multiple of CG's; the
\pmrep's MPI time ranges from 0.44 to 3.26 times CG's. The stand-ins are therefore useful,
where at all, only for choosing an arm, not for estimating CG's run time.

In answer to whether the benchmark is useful, ABCG's own SpMV exchange is the only
stand-in that met our 10\% criterion everywhere. The exact-matrix and per-rank replays
make the aggregate-or-not decision correctly when aggregation's benefit to CG is large,
and recommend against aggregating when the benefit is at most 1.21 times, because the
benchmark's LOCALITY path is 1.49 to 1.99 times slower than ABCG's. The aggregate-histogram
replay is not usable for this decision at 2 and 4 nodes. No stand-in chooses correctly
between MPI and STANDARD, though that choice cost CG at most 4.2\% in this data.

\section{What This Data Cannot Show}
\label{sec:limits}

This campaign leaves several questions open, and each limits the conclusions above.
First, each run reports only totals over $n$ iterations, so we have no per-iteration
times, no tail percentiles within a run, and no test of the plan's extreme-value
hypotheses. Second, the replay benchmark closes every exchange with a barrier and does
no computation, so it cannot show accumulation of delays across iterations or overlap of
communication with computation. Third, all data comes from one machine, one MPI library,
one MPI Advance version, and a shared queue over 23.7 hours; the spread across
allocations is likely smaller than it would be for allocations weeks apart, and nothing
here says the arm orderings hold on another machine. We also did not record Dane's
interconnect. Fourth, the campaign runs three arms rather than the plan's five, no
one-sided arm, and no combined Vernier manager. Fifth, CG never converged within 500
iterations on these matrices, so CG time here is time per fixed iteration, not time to
solution. Finally, we do not know what causes the 0.160\,s stalls, so we cannot say
whether they belong to Dane, to Open MPI 4.1.2, to these patterns, or to the up to 16
launches running at once. Timing individual iterations of the 4-node crankseg\_2 MPI runs
would show whether single iterations take about 160\,ms.

\section{Conclusion}
\label{sec:conclusion}

In this paper we measured how much ABCG's CG iteration slows down on Dane when each of
four stand-ins chooses among plain MPI, MPI Advance neighbor alltoallv, and MPI Advance
locality-aware aggregation, over 40 allocations at each of 2, 4, 8, and 16 nodes and six
matrices. ABCG's SpMV exchange chose within 4.2\% of CG's best everywhere. The replay of
the exact matrix halo and the replay of the per-rank Vernier capture chose within 10\% at
23 and 22 of 24 combinations. The exact-matrix replay missed at exactly the three
combinations where aggregation's benefit to CG was at most 1.21 times, because the
benchmark's locality-aware path is 1.49 to 1.99 times slower than ABCG's. The per-rank
replay matched the exact-matrix replay except at one near-tie. The replay of the
aggregate Vernier histograms chose differently from the per-rank replay at 4 of 24
combinations, all at 2 and 4 nodes, where CG took up to 32.3 times as long per iteration
on its choice. Every stand-in, including ABCG's SpMV exchange, chose STANDARD where CG
was fastest on MPI, at a cost to CG of 0.1 to 4.2\%. Finally, 266 of 9,600 MPI and
STANDARD runs lost whole multiples of 0.160\,s inside the timed loop. By median,
LOCALITY was CG's fastest arm on the crankseg matrices at 2 nodes, but on each arm's
fastest 10\% of launches MPI was, so these stalls decided that ordering. The aggregate
replay almost never stalled, and we have not identified the cause of the stalls.

\bibliographystyle{plain}
\bibliography{references}

\end{document}
